Ce projet de fin d'études m'a permis de réaliser \textbf{PharmacieConnect}, une plateforme web et mobile
liée à un besoin concret : faciliter l'accès aux informations sur les pharmacies et les pharmacies de garde,
tout en proposant un espace de gestion pour les administrateurs. La solution regroupe une API backend,
un tableau de bord administrateur et une application mobile.

Les objectifs principaux ont été réalisés. L'utilisateur mobile peut rechercher des pharmacies, utiliser la
géolocalisation, consulter les pharmacies de garde par date et accéder au catalogue de médicaments.
L'administrateur peut gérer les pharmacies, les gardes, les médicaments, les utilisateurs autorisés, les
imports CSV, les indicateurs et les journaux d'audit. Le backend assure la centralisation des données,
l'authentification, le contrôle des rôles, la validation des API et la traçabilité des actions sensibles.

Sur le plan technique, ce travail m'a permis de mettre en pratique plusieurs compétences : conception
d'API REST, développement backend avec FastAPI, modélisation SQLAlchemy, développement web avec
React et Vite, développement mobile avec Expo et React Native, authentification JWT, validation de
fichiers CSV, journalisation, analytics et tests automatisés. Il m'a aussi montré l'intérêt d'une séparation
claire entre les routes, les services, les modèles, les interfaces clientes et la base de données.

Certaines limites restent présentes. Le catalogue des médicaments ne représente pas encore le stock réel
de chaque pharmacie. Les données de garde dépendent toujours de la qualité des fichiers importés. Les
tests end-to-end et les tests de charge doivent encore être renforcés. L'application mobile devra aussi
être testée sur davantage d'appareils et dans différents contextes de connexion.

Les perspectives d'amélioration sont les suivantes :

\begin{itemize}
    \item ajouter la disponibilité des médicaments par pharmacie ;
    \item permettre aux pharmacies de mettre à jour directement certaines informations sous contrôle ;
    \item envoyer des notifications lorsque les gardes changent ;
    \item améliorer la recherche par synonymes, erreurs de frappe ou noms commerciaux proches ;
    \item intégrer des tests end-to-end couvrant l'application mobile, le dashboard et l'API ;
    \item renforcer le déploiement par une configuration de production, une supervision et des sauvegardes ;
    \item enrichir les analytics pour mieux suivre les recherches et la couverture territoriale.
\end{itemize}

Ainsi, PharmacieConnect représente une base fonctionnelle pour une solution de gestion pharmaceutique.
Le projet peut encore évoluer, mais il montre déjà comment une architecture bien organisée peut servir à
la fois les besoins du citoyen et ceux de l'administration.
